% ============================================================
%  F. C. Sibbern, On Cognition and Inquiry.
%    Toward an Introduction to Academic Study.
%  Copenhagen: At Fr. Brummer's Press (Printed by C. Graebe), 1822.
%  TRANSLATION (English) — SCAFFOLD.
%  Source of truth for translation: transcription.tex (Danish, COMPLETE + proofed).
%  Method: TRANSLATION-PLAYBOOK.md. Fill each "% [text to be added: ...]"
%    marker batch by batch (~10 printed pp.), mirroring transcription 1:1.
%  Register: lightly modernized (faithful meaning, syntax untangled).
%  Key terms: on first use in a section, gloss the Danish in brackets, e.g.
%    cognition [Erkjendelse]. Maintain GLOSSARY.md for consistency.
%  Page-offset (Danish scan, for reference): body printed = PDF - 22.
%  See TRANSLATION-NOTES.md for state + the section->page map.
% ============================================================
\documentclass[12pt,a4paper]{book}
\usepackage[utf8]{inputenc}
\usepackage[T1]{fontenc}
\usepackage{libertinus}
\usepackage{libertinust1math}
\usepackage{textalpha}   % real Greek glyphs, copied verbatim from the Danish
\usepackage[english]{babel}
\usepackage{enumitem}    % for \arabic*)-style lists (playbook: easy to forget)
\usepackage[protrusion=true,expansion=true]{microtype}
\usepackage{setspace}
\usepackage[margin=1.2in]{geometry}
\usepackage{fancyhdr}
\setlength{\headheight}{28pt}
\addtolength{\topmargin}{-13.5pt}
\usepackage{hyperref}

\hypersetup{
  pdftitle={On Cognition and Inquiry},
  pdfauthor={F. C. Sibbern},
  hidelinks
}

\pagestyle{fancy}
\fancyhf{}
\fancyhead[LE,RO]{\thepage}
\fancyhead[CE]{\textit{\leftmark}}
\fancyhead[CO]{\textit{\rightmark}}

\setstretch{1.3}

% TITLE — rendering to confirm with Hans (see TRANSLATION-NOTES): "Granskning"
% could also be Investigation / Scrutiny / Examination.
\title{\textbf{On}\\[6pt]
  \Huge\textbf{Cognition}\\[6pt]
  \large and\\[6pt]
  \Large\textbf{Inquiry.}\\[12pt]
  \normalsize Toward an Introduction to Academic Study.\\[16pt]
  \normalsize By\\[6pt]
  \large Dr.\ Frederik Christian Sibbern,\\[4pt]
  \normalsize Professor of Philosophy.}
\author{}
\date{Copenhagen 1822.\\
  At Fr.\ Brummer's Press.\\
  (Printed by C.\ Gr\ae be.)}

\begin{document}
\maketitle

\frontmatter

\chapter*{Preface.}
\markboth{Preface}{Preface}

\begin{center}\rule{0.28\linewidth}{0.4pt}\end{center}

There are many writings on the study of the sciences [Videnskaber] in
general, and on academic study and the student's vocation [Kald] in
particular. The present one is a new attempt of this kind, following a plan of
its own. If giving oneself a full account of what one does and accomplishes
can do nothing but advance a person in his striving and activity---so that he
may come to act with complete clarity and composure---then this must be
especially important for the one whose whole endeavor is precisely to examine,
everywhere, the innermost nature and essence [Væsen] of things. And though
there may admittedly be, in other spheres of activity, people who, instead of
being advanced, would be held back and arrested---indeed led astray in their
striving---were they to strain to penetrate into their own selves and search
out all that underlies their active power and the natural gift they have to
thank for it, this cannot hold of the student as such; for he, called to
penetrate to the ultimate grounds of all cognition [Erkjendelse] and of all
things, and above all to ponder the nature of spirit [Aand] and of the
spiritual life, cannot exclude from that inquiry the form of life that emerges
in his own inner being.

This is why, in our own time, reflections on academic study have been made the
subject of separate lectures at the universities.

But if one considers more closely what belongs to such a treatment and what
does not, difficulties arise---difficulties that the author of the present
attempt, too, could not but feel. To clarify the nature and essence of
scientific study, one must enter into considerations that other scientific
disciplines nonetheless claim as their own, and the question then becomes how a
boundary is to be drawn between the two. Let me be permitted a few words on
this matter here, so as to clarify the relation in which this presentation is
meant to stand to the other scientific developments it meets at several points.

To show what a person who strives after genuine cognition and science
[Videnskab] ought to keep in view, the first and most important thing is to
consider what cognition and science are. But underlying all cognizing lies
thinking [Tænkning]. The ordinary developments of logic, however, cannot belong
here. Just as little does it belong here---though the life that is to live and
take shape in the scientific inquirer [Gransker] is indeed the object of our
consideration---to set out the various life-activities that are to be
distinguished in our spiritual life, together with their relations and
connections, or to enter into psychological investigations. Nor should an
introduction to academic study be confused with that general introduction to
philosophy which has, with good reason, lately been established as a discipline
of its own. Yet this latter cannot but come into manifold contact with all
these investigations. And so the question arises what an introduction to
academic study really is to be and to accomplish; indeed, doubt [Tvivl] may
well arise whether it can even claim and fill a place alongside those other
scientific presentations.

We bring this matter to clarity most easily by asking how we came---how others
before us have come---to make academic study a distinct object of distinct
investigation. It was because it had to matter to the student, who as such has
his own particular vocation, to reflect on this vocation of his, on his
striving and activity. Among the forms in which the general spiritual life that
psychology considers can come forward in a particular individual existence,
hypostatizing itself in a distinctive way, the life that is to live and dwell
in the student is one; in other individuals, in other spheres of activity, it
takes shape in other ways. What now serves to illuminate this form of life from
every side, and to remind us what it really is that we as students strive for
and ought to strive for---and what we, in the midst of a life in which so much
can distract and confuse, or draw us toward one-sidedness [Eensidighed], hold
us back and detain us, ought constantly to keep in view---this is what belongs
here. And thus there appears a distinctive center around which the manifold
[Mangfoldighed] that presents itself to reflection must gather into unity
[Eenhed].

A presentation of this kind may be regarded as a sort of monograph. Of such
self-standing scientific treatments of the whole nature and character of
certain particular objects---or of objects belonging to larger scientific
systems---one can conceive two kinds. The usual natural-historical ones, and
those resembling them, treat their object exactly as it is otherwise treated in
its place within the doctrinal edifice to which it belongs, and so are to be
regarded merely as excerpted sections or pieces of the larger scientific
whole---perhaps only elaborated more fully than usual. Of a different character,
by contrast, is a monograph such as the present one. Items of knowledge that
otherwise belong in the most diverse places in other sciences are here to
gather about a single point, in order---together with, and in one with, the
reflections on what is peculiar to the object---to illuminate it fully and from
every side. Such a monograph, more difficult and intricate than one of the
first kind, also has a better claim to stand forth as a self-subsisting
scientific inquiry, or as a smaller whole within the great whole of the
sciences; for here there is a distinctive center for a scientific organization
of its own, whereby a body of interconnected cognitions grows together into
unity in a peculiar way. And this organization---the inner ground-form of the
presentation, the inner mode of connection and sequence of its parts, the
course of its development---is determined by the object's own distinctive
fundamental nature, not by a norm and manner of presentation borrowed from
another science and fixed by that science's course and method [Methode]. A
monograph of this kind stands to the sciences to whose fundamental
considerations it attaches itself---and to which much of what it presents and
works out is to be traced back, though always in a distinctive way determined
by the object, and in a light of its own---as a talent stands to the general
activities of the mind. Take away from a talent all that belongs to
understanding [Forstand] and judgment, sense and imagination, reason [Fornuft]
and feeling, and nothing will remain that could be called a talent; and yet a
talent is something other than those more general activities of the mind. What,
then, is a talent? It is the distinctive combination of them all in certain
definite directions---the distinctive organization of the mind with respect to
their united cooperation, aimed at a certain kind of productivity. Just as a
thorough psychological portrait of a given talent does not consist in a
portrayal of those more general mental activities in general, but must
presuppose them as known---while it is precisely they that it must everywhere
attend to and keep in view---so a monograph like the present one may have to
look toward many particular points which, taken in and for themselves, belong
to definite places in definite sciences, and may have to present much that,
piece by piece, other sciences could claim and draw to themselves, and can
nonetheless be, and assert itself as, a distinctive scientific whole in its own
right.

It was, moreover, my intention from the outset---the reader will perhaps notice
this easily in several places---to lay the content of the present work before
the public in a form other than the present one: in the same form in which I
have several times imparted it to my listeners, the form of lectures. The
lecture is, among us, without doubt the mode of presentation that could most
easily take the place the dialogue holds in Plato; like the dialogue in his
day, it is taken directly from life, and the development of general concepts can
hardly, in any way natural to us, be better combined in a connected
presentation with speech's immediate address to individuals---and take on the
liveliness that must follow from such direct attention to the thinking beings
who are to be called upon to meet what is presented with their own independent
thinking [Selvtænkning]---than in the form of the lecture. But the fear that I
might not come even remotely close, that my attempt might fall all too far short
of the ideal of such a presentation which hovered before me, brought me to
abandon the execution I had already begun.

I had also intended from the outset, in the Remarks to \S\S~13--17, to treat
the individual principal fields of study, both with respect to the philosophical
and in the other principal respects. But what I might have said about this I
later decided to set aside and reserve for the oral lectures I hope still to
give on academic study, for those students who are still at the beginning of
their course. I dared not credit myself with such an acquaintance with so many
sciences that I should be able, even in general terms, to offer reflections on
their study fit to be laid before the public in print, in a word fixed for
good---or at least for a long time.

I must further note that if, in the present work, I have departed from the
practice I observed in the two scholarly works I have previously
produced---namely, of emphasizing no word---this has happened because I believed
I ought to yield to the urging of one of those among my friends to whose company
I owe much, especially with respect to my philosophizing endeavor. But since I
am not accustomed to underline as I write---because it almost never occurs to me
that it might be necessary---and since this time I had nearly finished the whole
book before the decision to do so was prompted in me, so that I carried out the
underlining only at the work's final revision, whereas the moment of writing
something down is the right moment to decide which word should be
emphasized---I dare not hope that I have carried out the decision with the best
success, and I therefore ask the reader, who will in any case find enough else
to overlook, to extend his forbearance to this too.

This is what I had to send ahead of the work I herewith deliver into the hands
of the public.

\begin{flushleft}
Copenhagen, 13 November 1821.
\end{flushleft}
\begin{flushright}
Sibbern.
\end{flushright}

\begin{center}\rule{0.28\linewidth}{0.4pt}\end{center}

\chapter*{Synopsis.}
\markboth{Synopsis}{Synopsis}

\begin{center}\rule{0.28\linewidth}{0.4pt}\end{center}

All cognizing [Erkjenden] is essentially a finding, a receiving, and a
beholding. Truth, presenting itself out there of its own accord, refers itself
to an apprehending individuality and carries with it, for the latter, a
necessity that produces conviction [Overbeviisning], to which a personal
affection is joined. (\S~1.)

Yet we cognize only insofar as we ourselves bring forth. \textit{Scimus, qvia
facimus.} A living personality cannot receive except in and through its own
forth-working. (\S~2.)

Receiving and producing here become one. Precisely by ourselves shaping truth's
expression or presentation within us do we receive it and are made alive by it.
(\S~3.)

In pure beholding, however, productivity is lost and does not come into view.
But in the striving to cognize perfectly, reflection makes it become free and
show itself as such. If we ask why we come to truth only slowly, through
reflection---why its fullness is not given at once---the answer is: because a
realm of freedom is to come forth, so that a realm of love may become actual.
(\S~4.)

The thinking called forth by reflection now leads to a cognizing by force of
grounds [Grunde]. The highest cognizing would indeed be a pure beholding of
perfect truth, raised above all inquiry [Gransken] and grounding; but such an
immediate beholding of the highest is not granted to the human being. Rather,
all his assurance rests, in the end, on an immediate faith [Tro]. (\S~5.)

The endeavor to cognize by force of grounds is an endeavor to cognize truth in
its truth. This is what calls forth genuine scientific study, which first truly
begins with academic study. (\S~6.)

Cognition now becomes an intelligent cognizing (as opposed to the empirical); a
cognition of truth in its truth is necessarily a cognizing by force of the
constitutive First, hence a cognizing \textit{a priori}, which leads to a
construction \textit{a priori}. (\S~7.)

Such a cognizing then aims not merely at grounds for our assuming something, but
at things' own grounds of existence [Tilværelsesgrunde]. All philosophizing
strives to find these. (In mathematics, by contrast, we have a great example of
a cognizing---even a wholly intelligent one---that does not aim at this.)
(\S~8.)

The real grounds of things we find, most immediately, in the rational
orderliness [Forstandighed] that, within a coherent order of things, requires or
entails them---indeed, in the concepts [Begreber] of things. (\S~9.)

But as the ultimate ground there appears an all-real active First which, in
bringing a determinate order of things into actuality, sets the individual
things into existence in accordance with a rational lawfulness
[Forstandsmæssighed]---indeed, in accordance with concept. (\S~10.)

This all-real presents itself to us in the Idea [Idee]. Cognition is in itself
the Idea's own self-presentation in consciousness, even if it does not always
show itself---or is not seen and possessed by us---as such. That truth is not a
mere correspondence between an assumption and its object, but is that which
livingly produces this correspondence, is clear from this. (\S~11.)

The Idea presents itself to cognition and comes to outward life and subsistence
in science. (\S~12.)

The first thing in all genuine science is the Idea itself, as science's center
and all-determining soul. In the cognition of it, and in the tracing-back of
all cognitions to it, the speculative and philosophical element in science shows
itself. Yet philosophy can also show itself already in the immediately
judicious---in a certain scientific way of thinking and judging. (\S~13.)

In speculation's transition to the manifold---which is thereby to be beheld in
its proper light---and in the determination of that manifold, dialectic enters,
partly as a philosophical art of annihilation, partly as a scientific art of
construction essentially and immediately bound up with it. (\S~14.)

Genuine cognition is, further, necessarily and essentially systematic. Yet the
immediately free judicious element, just as it must prepare the way for the
system, must also not be displaced by it. (\S~15.)

Furthermore, in all science the exact study of the particular [Detail] is of
importance. (\S~16.)

Finally, in all science method [Methode] is still to be noted. It is best
developed out of the science itself and its historical course of
development---just as the study of a science's history is of essential
importance for the study of the science itself. (\S~17.)

Not only in science, but also, through science, in a personal individuality does
the Idea take shape, and thereby come to a complete outward existence or
subsistence in real actuality. (\S~18.)

To attain this subsistence and life in the individual, the Idea then calls forth
its own opposite, setting the individual's sense of independence in motion and
summoning a recognition and avowal [Er- og Vedkjendelse] that is to follow on
the individual's part. (\S~19.)

Independent thinking [Selvtænkning] is called forth by a questioning. The
possibility of assuming, but also of not assuming, then appears. Denial and
affirmation stand opposed to one another. But the assurance that arises out of
doubt [Tvivl] beaten down---the Yes that comes forth from the denial of a
negation---is of an entirely different nature than what has not passed through
this trial. (\S~20.)

The Idea's coming to the individual then appears as a kind of testing, and
everything depends on whether that now follows in the individual which is to
follow. There is indeed something here that lies within the individual's power;
but that on which everything depends does not lie within his power, and must,
when it is given, be received as pure grace. The conflict that is here to be
brought to unity is seen above all in theological study. Yet in several respects
there is a trial to be undergone here, since, on the one hand, egoism easily
comes to draw nourishment from the consciousness of the independent thinking to
which the summons has been made, and, on the other, the very endeavor after
insight can easily lead into uncertainty, confusion, and error. (\S~21.)

Finally, the Idea will live and take shape not only in individuals, but in a
whole of humanity, in which each is to find himself as a member and to intervene
practically in the exercise of a definite vocation [Kald]. In the choice of that
vocation, science and Idea must themselves guide us. We must, further, take
science with us into life, and the exercise of one's vocation must at the same
time be a constantly continued study. (\S~22.)

But the highest, the one thing needful---however lofty and significant genuine
cognizing may be---is nonetheless not yet cognition itself. (Conclusion.)

\begin{center}\rule{0.28\linewidth}{0.4pt}\end{center}

\mainmatter


\begin{center}\large\S.~1.\end{center}
\markboth{\S~1}{\S~1}
All cognizing [Erkjenden] is essentially a finding, a receiving, and a
beholding. (\textit{Scimus, qvia accepimus.})

For all cognizing aims at the true, and in and through it we are to come to
truth [Sandhed]---to come to have and possess it. But truth is not something
dependent on us; it is and remains, without any regard to us, what it is in
itself. Ours, then, can only be to seek quite directly to become conscious of
it as it is in itself and presents itself to us; but such a coming-to-consciousness
is a receiving, and---when it is determinate and clear and at the same time
immediate, or certain and self-evident by itself alone---a beholding. But this
immediate, self-certain element is present at every point in all cognizing,
even when we reach it in a mediated way; for cognition then still always rests
on the insight that, from such prior cognitions, the present one follows---and
this following, or consequence, must be self-evident directly, or must be
beheld. With this, what thus follows likewise shows itself as that which, on
the presupposition of the truth of the prior cognitions, is self-evident,
though only on that same presupposition is it also cognized. So that even
relative and mediated cognition, at the point where it first enters by virtue
of its mediateness, proceeds from a direct insight, and so far is itself an
insight and a beholding. It is a beholding that gradually develops. In every
case it is a receiving.

As such, it is always and essentially bound up with a necessity to represent
[forestille] just as one does represent, and not otherwise. The immediate
feeling of this necessity is conviction [Overbeviisning], which can arise only
in that we are immediately seized by it and feel it. The object of cognition
presses itself upon us, and precisely as it is now cognized to be.

The true that is cognized thereby shows itself as that which exerts a
determining and effective influence in us---an influence in which it expresses
and presents itself, and thereby lives again, in us.

\begin{center}Remark and Supplement.\end{center}
1.~That all cognizing is a receiving and a beholding is soon taught, too, by a
glance at the various domains of cognition. Not only what strikes the eye at
once: that the \emph{beginning} of all cognition---in the higher spheres just
as much as in the lowest sensory one---is a direct becoming-aware and
apprehending, a noticing and sensing. In its advance toward the highest
completion, too, it remains an immediate taking and beholding. In the domain of
lower sensation this is at once, at first glance, entirely obvious; for although
here a certain combining and synthesizing activity must indeed be added to the
apprehending one---so that sensation comes about only by means of a reproducing
construction---this construction too aims only at presenting the thing exactly
as it is found and as it presents itself to us. We feel ourselves affected, and
just as the affection brings it about, so sensation turns out: precisely as we
are determined to it by this affection, so must we represent. And though a
difference not seldom arises here between what seems to be and what is found to
be in itself, still this latter too is found, and must precisely be cognized,
only as it is found through a scrutiny---if not by sense, then by deliberating
reflection.

As it here shows itself in the sphere of the senses and of the experience they
occasion, so it stands with every higher cognizing of the understanding
[Forstand] or of reason [Fornuft]---indeed, it is especially striking in the
highest cognizing. Whatever is cognized without appeal to experience, by an
insight springing from pure thinking---whether by force of mathematical
demonstration or of a priori reflection or of philosophical speculation---is
always cognized as something that is seen to have to be so, and a necessity of
understanding or of reason, determining our whole representing, thereby enters:
as this necessity declares it to be, so it must be cognized to be. But when, in
this way, all reflection---all comparing, connecting, and inferring---aims at
investigating or fathoming what must necessarily be assumed, as it will itself
turn out to be, so that it is cognized as true only in being cognized as the
one thing that can be assumed, then reflection aims at letting the thing
\emph{give} itself, by letting it be forced upon one through a necessity of
thinking. It is cognized only in being discerned---that is, seen and beheld
(with the higher eye of the spirit [Aand])---as it is in itself. The triangle I
construct; but that its angles together equal two right angles I discern as
necessary in accordance with the triangle's own nature, once given to me.
Indeed, in the very act of constructing it, I do so in a certain determinate way
and under certain conditions, and this manner of construction and these
conditions I must take as they force themselves upon me. All sorts of human
circumstances and relations I can think out for myself (though even here I must
proceed according to certain given conditions and determinations); but what, in
a given relation, is right or wrong, good or bad, the fundamental Idea
[Grundidee] of the good and the right must teach me, so that by it I must let
myself be guided in my judgment, and in every case let stand as right and good
what, in accordance with that Idea, I see to be right and good. Mine it can only
be to inquire into the essence [Væsen] of the right and the good, and then to
cognize it as it presents itself to my spiritual eye. So that, even if the
cognizing which goes beyond mere experience is called forth by questions---which,
as inner spurs to seek more than the experienced, show that we do not rest
content with what is given in experience alone---still the whole thinking
thereby set in motion leads only to a new receiving and apprehending of a
higher order, under which we must again seek directly to observe and behold what
presents itself to us.

But now, in the receiving and beholding in which all cognizing consists, we also
become aware of a \emph{principal distinction}. What we apprehend and come to
know through the senses---indeed also what the ideal or intellectual
becoming-aware and beholding, which enters already in this sphere by way of
thinking, can teach us, when the inner ground-forms of the things revealed to
the senses, their relations and their activities, are cognized exactly as they
are directly found---all this is merely received, as something given to us from
without in the reality that presents itself in time and space. This cognizing is
the merely \emph{empirical} kind; the necessity by which it enters, and by which
it acquires the character of a cognizing, is only a necessity of the sort that
proceeds from a certain \emph{affection} bound to time and space---a wholly
sensory affection of a lower nature. What is cognized in this way is still
cognized only as \emph{contingent}. It is given to us as the actual, and for us
there is a necessity to take it so; but that it should itself necessarily be as
it is, is not yet seen in this merely empirical cognizing.

Of a wholly different kind, by contrast, is mathematical and philosophical
cognizing; here there is an intellectual beholding that has nothing of that
stamp of contingency. What is here cognized is not merely taken as given, as
though that were enough, but is apprehended and represented as something that is
so because it must be so, in accordance with an inner \emph{necessity of
understanding and reason} of a wholly different origin than that by which
certain sensory affections force themselves upon us. It is indeed given to us,
but given from within---or at least given as something that accords with an
inner demand, and not at all given to us by way of sensory affections. It is
cognized as something that cannot be otherwise; it is cognized as
\emph{essential}, and therefore valid for all times and everywhere---which is
why only in this sphere can one speak of eternal truths. We will call this
cognizing the \emph{intelligent} kind, both because this word seems more apt
than the word ``intellectual,'' and because the latter can also cover the
just-mentioned merely empirical---yet still ideal---cognizing, in which the
inner ground-forms and relations of things are apprehended entirely as they are
directly found in experience.

2.~Necessity enters when truth shows itself in relation to (refers itself to)
\emph{the individuality} that takes it up; and in the feeling of conviction, the
life and working of the true in the individual is met by something purely
individual, in which a person shows himself and comes forward as such, in that
he attaches himself wholly to, and gives himself over to, that working. In pure
receiving and beholding, the true was as yet seen only as presenting itself, so
that only its own (universal) appearance showed itself therein---no personal
life. The reception was merely like a reflection or mirroring, like that by
which water gives back light and color---not a personal taking-up, like that by
which the gemstone has, in its color, taken a power and effect of light into
itself, into unity with its own nature and existence. But in conviction it
becomes an inward appropriation, whereby the true that presents itself in and
for an individuality now also comes to dwell powerfully in this individuality,
as in a being in which it has taken its seat, and which itself now is and
becomes something through the reception. In this binding-together of the true
and the individual into unity, the necessity that is indeed present is so lost
that it is not particularly felt or brought into view. Here there is an
\emph{undivided} free life in the true, raised above the feeling of necessity.
Only when the individual, tearing itself loose from the universal, moves in a
particular personal striving does the opposition---and with it that same
necessity---come to light.

3.~In conviction, a personal \emph{affection} and emotion, πάθος, is added to
the reception---with a personal attachment and devotion to the power of
truth---so that the individual not only lets truth gain room and power within
him, but does so with delight, feeling that in this he possesses something and
that his personal life is thereby expanded and heightened. The devotion now
shows itself as \emph{love}. This brings it about that, together with the power
of the true, personality and freedom too come fully forward; for in love an
independent being goes willingly and gladly to meet the true. Whereby the
endeavor to cognize shows itself as a cultus or \emph{worship} of truth. The
inquirer's [Gransker] relation to truth and to his science [Videnskab] now
becomes a \emph{religious} relation---in the not ordinary, but literal and
far-reaching sense in which the word denotes a relation, bound by sacred ties,
to something holy, in whose power we feel ourselves, with willing devotion, to
be. (Indeed, in the word's more ordinary and proper sense too the relation can
show itself as religious; for just as a human being can and should do all that
he does with God in view and out of love for God---so that this highest motive
can and should be the everywhere-operative one---so too not only can that same
motive, as such, embrace the scientific endeavor as well, but everything can and
should in the end be cognized and beheld in its relation to the highest divine
in such a way that the Idea of God becomes, scientifically, the all-illuminating,
just as, practically, the all-animating. So that in the end every relation to
truth shows itself directly as a relation to God. But in this we look toward an
ideal that the scientific culture [Videnskabelighed] of humanity is still far
from having approached.) With reverence for the voice of truth, and a heartfelt
desire to bring it---and it alone---to speak and live in his inquiry, the votary
of science must give himself over to being its organ, its servant and herald,
who wills only what it wills. His personality must be as a holy temple for it.
But at the same time the speaking and living of the true within him must be felt
and held fast as his own innermost and best life, and the relation must show
itself as a true relation of love. What has sometimes been said---that a person
is consecrated to his science and his art---we may recall here. To be kindled by
truth, and then to take it up and faithfully preserve and tend it within us,
until it grows and develops of itself and takes shape in us, so that it may
truly come to fulfillment in us---this is ours. Let love of truth
[Sandhedskjerlighed] be, like all genuine love, careful and conscientious.

Truth thereby shows itself as that which, \emph{by its own might and power},
seeks to take shape in the individual, while the individual gives himself over
entirely to letting it work in him whatever it will. And just as this holds of
everything that, in any science or body of knowledge, shows itself as the true
whose cognition study aims at, so everything set forth here holds of every
scientific endeavor---even of one that still stands so far from the higher Ideas
that it only rarely feels itself inspired by them. In \emph{every} study,
faithful devotion to the cognition that forms itself in us will bring the
fullness that love gives. If, in the lower domains of inquiry, what we cognize
seems not to bring us nearer to the eternal, still the fact that it produces
complete cognition and living conviction in us will make us feel as though
filled by a divinity. (Just as, even in the lowest spheres, cognition in itself
really does come from the higher source, even if it does not come to
consciousness as such. See \S~11, Remark~2, below.) So that every study in which
faithful and clear inquiry prevails has its \emph{Muse}. In it the individual
stands over against something he would make his own, needing it as something
without which he himself cannot come to full life. But what he would thus make
his own is not something he can treat as an object for the satisfaction of a
sensory need---to consume, or to take as a mere vehicle or means. It always
remains standing over against the individual as something that subsists in and
for itself in its significance and worth, whether or not it comes to the
individual---indeed, as something that has worth for the individual precisely
only insofar as it has worth and significance (that is, in this case, has truth)
in and for itself. And thus, everywhere that truth is sought, in whatever
sphere, the relation between the individual and the true he seeks is such a
relation of attachment as can---and, under a genuine striving, also will---become
a relation of love.

We now also see here what it is to \emph{inquire}, to \emph{fathom}, to
\emph{reflect}---or, as one less felicitously puts it, to carry out reflection,
investigation, consideration upon something. Namely, that it can be nothing
other than, under this or that arrangement and handling of certain
representations, to try to bring it about that the truth we desire to know may
come, of itself, by its own might and influence, to present itself to us---so
that we now directly become aware of it or, as one says, find it out.

We must therefore also \emph{immerse} ourselves in every study. Only thereby
does it become genuine and powerful. If the true that produces cognition in us
is to gain us into its power, so that it may form and shape itself in us, then
we must forget everything else over it and, as it were, lose ourselves in it.
Even the sensory forms of things are only then observed and apprehended fully in
their distinctiveness and manifoldness when attention is thereby wholly
captivated, as though we were entirely absorbed in the contemplation, or lost in
it; for only then does the living \textit{nisus formativus} that here
enters---under which the true would take shape in us---work whole and undivided
in its full strength. A book is best studied and works most fully when, for a
time, it so captivates us and, as it were, takes our whole mind into possession,
that we are wholly immersed in it. In this way---that is, with this vigorous
attention and devotion---should every \emph{classical} work be studied.

But only \emph{classical works}, too, deserve and can be wholly studied thus.
For a scientific work is classical not by the collection of items of knowledge
it may contain, and which can now be drawn and learned from it piece by piece,
but by its spirit as a whole, its inner genius, which speaks out of it in its
totality, and by the total form that accompanies this---the way in which
everything in it is presented and composed and, as it were, grown together into
unity, so that everywhere in the work we meet the same vigorous and profound
spirit, and all the manifold in it is as though of one piece. Such a work
stands, by its spirit and by its form, valid and significant for all times,
which the word ``classical'' especially seems to indicate; for it works not
merely by a manifold of material that can be transplanted and reworked, but
works as a whole by its distinctive original stamp, whereby spirit speaks to
spirit. But then this spirit too must work in its totality and, as it were,
enter into the reader or beholder---and it does so all the more, the more he
immerses himself in the contemplation, so that he treats the work not merely as
a means to the acquisition of knowledge, but as a work that has full worth by
and in itself alone, and therefore also works toward the very formation
[Dannelse] (\textit{Bildung}) of the spirit as a whole.

We must read in such a way that we enter into a \emph{harmonious relation of
reciprocal action} with what we read, and are so set in motion that, through the
words and sentences, we come to feel ourselves addressed by the general
spiritual element in the work, which constitutes the inner source and first
mover (\textit{primus motor}) of the words and sentences---the soul that lives
and speaks in them. Only such reading is a truly soulful reading. The study of a
work thus becomes an inner communion with a glorious spirit, which in several
respects can have a beneficial influence on us and advance us in our whole
striving. It becomes, throughout, all the more genuine nourishment for us the
more it is a \emph{disinterested} study; for the more we inquire and probe not
for the sake of this or that particular aim, but solely out of pure love of
truth, in order to take it up entirely into ourselves, the more it will work in
us with its full power. The more our devotion is undivided, free, and
wholehearted, the more easily and fully can the total immersion, or
being-gripped, come about. For a time, then, we let the work carry us away, and
let it do with us, and make of us, what it will. Thereupon truth's own power
itself begins to speak in us and to bear witness to what is expressed in the
work; whereby an immersion in our own immediate contemplation of the
intrinsically true is added to the immersion in the work's utterance, and now
determines what the work is to make of us.

In general, the spirit that lives and addresses us in a work can, of course,
come to carry us away and captivate us, and become for us a source of new
effects and formations of life within us, only in that it comes from truth
itself---indeed, from that which can present itself immediately to every
inquirer, enter into everyone, and in everyone bear witness to itself. Only in
that we find the same thing again in ourselves, so that we sympathize with it,
can what is handed down to us speak to us through the forms of transmission in
which it is imparted. (This holds even of the study of the works of nature;
nature too is a book that must be studied like any other book, but in which a
spirit akin to us also speaks to us: the same truth that bears witness to itself
within us as well. Yet this will only become fully evident in what follows. See
\S~7, in the Remark.) To grasp, cognize, and enjoy a work handed down to us we
come, in the proper sense of the word, only in that it \emph{reproduces itself
in us}, and we, in the course of this its reproduction, are moved and enlivened
by the spirit which---just as it now lives again in us---first lived in the
activity by which the work came into being. But by this we are led directly to
the next principal consideration, as we take into closer examination this
reproduction, which shows itself as one whereby we ourselves are productive.

\begin{center}\large\S.~2.\end{center}
\markboth{\S~2}{\S~2}
All cognizing comes to actuality and subsistence by means of a determining, a
doing, a producing---and only by means of such; so that the latter is
essentially bound up with the former. \emph{We cognize only insofar as we
produce.} (\textit{Scimus, qvia facimus.})

Cognition takes place when the true shows itself as such to a personality and is
taken up by it. But a personality must reveal itself in a personal life. Now a
life cannot be conceived without a Something, a real [Reale], in which it
expresses itself in uninterrupted activity, continually forming and determining
this real, or outward, of its own; it gives itself a definite shape in giving
this a definite shape. Personality, on the other hand, shows itself in finding
and positing itself as the immediate center of its own existence [Tilvær], and
in having this existence as its own. A personal life, then, must express itself
in an activity and productivity in which the personal being finds and moves
itself as the producing agent. If, then, the true is to be taken up into the
personality, it too must be embraced by this personal productivity. The
individual can therefore cognize only in that it itself works and produces a
real---a (psychic) outward---corresponding to the cognition, in which it holds
the cognition fast by determining and developing this outward.

Consciousness is the personality's first expression. In it is taken up
everything that is to be taken up into the personality. Thus we too can cognize
only by becoming conscious of what is to be cognized. But this, according to the
fundamental nature of consciousness, we can do only by presenting it to
ourselves, or representing it to ourselves [forestille] (objectifying it for
ourselves). But this presenting or representing is our own act of fixing,
whereby we posit (or \textit{ponere}) the given and received before us and, as
it were, make ourselves an inner image of it, in which we express it for
ourselves and hold it fast. In this representing consists, here, that
productivity. The representation [Forestilling] is our own product and, as an
expression of our own life which we keep as ours, cannot be anything else. But
only in this constant representing, and under its constant determination,
development, and completion, do we come to cognize, and do we have what we
cognize.

The possession of cognition therefore also shows itself as an \emph{ability}, a
\emph{capacity}. We know only (sufficiently and fully) what we are able to bring
forth in one form or another and thereby give to others. This productivity of
ours, by which the received is now embraced, is also in itself a free one, which
can arbitrarily present something other than the true. In cognition, however,
this free element is bound and determined to a single way of representing, by the
necessity that accompanies cognition and now comes forward in this respect too.

\begin{center}Remark and Supplement.\end{center}
1.~Here we see the essential connection of \emph{language and speech} with
cognition. All cognizing that does not aim at immediate or reproduced
sense-intuitions, or at constructions carried out directly in time and space in
the same way as these (such as, for example, the mathematical ones), takes place
by means of a constant judging, which in turn takes place through an inner
speaking, by means of words (or possibly other signs), and cannot come about in
any other way. But that this is so---which experience teaches directly as a
fact---also shows how immediately a positing, determining, and producing
activity is added to our apprehending and receiving. In word and speech we set
what is cognized before us, thereby binding it and taking possession of it,
making it something of our own, which we now, in our own way---in our
designating and naming---hold fast (cf. Hegel's \textit{Encyklopädie}, § 379 and
the following §§).

2.~Yet not only by words: in all representing generally we do the same. The
preservation of representations after the immediate beholding is past, as well
as everything we can do with them, testifies to this generally. The productivity
essentially belonging to all cognizing also soon shows itself, to the observing
eye, in every one of cognition's domains. As for simple \emph{lower sensation},
we refer to the psychological observations that lead directly to the insight
that spatial forms and relations, as well as temporal determinations and the
qualities ascribed to things in respect of these, are cognized only by means of
a reproducing construction of the objects that present themselves---observations
which, already merely as such, lead to conclusions that justify assuming the same
as regards the merely visible, audible, and other merely sensory qualities. (See
the author's \textit{Psychologie}, § 55.) How \emph{thinking} takes place, and
the cognizing springing from it comes about, by means of a constant determining,
positing, and constructing, the logical investigations teach throughout and
everywhere. Whether concrete objects are conceived and their nature thereby
cognized by force of concepts [Begreber], or concepts themselves are regarded as
the objects to be cognized, in either case construction takes place. For the
former happens only when a connection is discerned; but this connection is seen
only when the manifold whose connection it is is set apart from one another and
then joined again into a whole, by means of a reproducing construction or a
genetic procedure, under which we let the thing arise before us---that is,
ourselves form it for us in an image---and now also see everything in its place,
in that we see the necessity of setting it thus into the whole. Connection, and
the real significance of objects within a certain coherent whole, can be
cognized only by projecting and completing the intuition [Anskuelse]
corresponding to it; to cognize it is to produce, or at least to be able to
produce, its concrete expression. The concepts themselves, next---which are
thereby abstracted, as the connections thus beheld in things and the
ground-forms and real meanings thereby determined, or as the fundamental
thoughts expressing themselves in things themselves---are cognized only when
they are lifted out on their own, by cognizing their so-called content: that is,
the sum of determinations that, according to the concepts, must be found in the
object and can be asserted of it as predicates. Whereby the judging mentioned
above enters, as a determining, a fixing, and a pronouncing. Deductions of
concepts---whereby the fundamental thought that determines the concept's whole
validity and content is derived, or seen in its origin from a higher
cognition---take place, like inferences and proofs, through trains of thought in
which thoughts are linked to thoughts by determinations, whereby a whole of
thought is developed. The deduction of concepts takes place only through a
construction of the scientific whole, in which the concept now shows itself as an
essential member; and proofs likewise come about only through combinations. An
Idea [Idee], finally, is likewise cognized only in that it is either beheld as
the expression of a higher Idea---in which respect, on the side of form, it shows
itself as a concept---or else beheld as the essence, the determining being, and
moreover as the real principle, giving essentiality and existence, in real
objects and ground-forms; which objects and ground-forms are then not only, like
all others, cognized solely by means of constructions, but can also be seen as
the expression of the Idea only in being seen as those which through it become
what they are---something that can be seen only under a reproducing, genetic
presentation.

\begin{center}\large\S.~3.\end{center}
\markboth{\S~3}{\S~3}
Reception and production, beholding and one's own determining and doing, here
become entirely one, so that cognition, precisely in entering under the one
form, essentially also takes on the other. (\textit{Scientia et potentia in idem
coincidunt.})

In cognition, the production and presentation is not arbitrary but one
determined with necessity, and what determines it is precisely our receiving and
beholding, and the true that presents itself in and for it. As thereby the whole
manner and character of the representation [Forestilling] is determined, the true
expresses itself---but through our productivity, which it governs and determines,
and in our product, which now nonetheless becomes the expression of that true,
and, as it were, becomes the material in which the true realizes and
individualizes itself in us, while at the same time our own life too realizes and
substantializes itself therein. We are determined in beholding, but determined to
produce something ourselves---namely, to form representations that become ours,
yet such that we produce them precisely so as to correspond to our
having-been-determined. And thus we are determining and determined in one. What
is given to us or presents itself to us comes to clarity in us---indeed, comes to
be ours at all---only in what we ourselves do. And yet here too cognition takes
place only by means of a constant receiving and beholding; for that by which we
are determined we have only in this its determining influence, and the
representation likewise, as necessary in accordance with that having-been-determined,
takes on the same character. \emph{We produce as we see that we must
produce.}

The personality, insofar as it finds itself as a consciousness and
self-determination---or we, who call ourselves ``we''---find ourselves here
between \emph{two regions}. On the one side, the influence by which we are
determined, and to which we turn and look as we receive and are determined; on
the other side, our sphere of representations, which we ourselves work over and
make into what it becomes, in that we seek therein to hold fast the received and
make it our own. The former, as determining, above us, before us; as that which
is to be expressed and to come to shape in us---the more inward and active, that
by which we are governed. The latter, as our product, below us; as that in which
we express our life---the more outward, the psychic material of our life, by its
nature an object of our governing. Here there is always, on the one side,
something that is obtained and possessed merely through the mastery it exercises
in us, or the determining influence it has on our working and determining; on the
other side, something that is merely expression, presentation, or result, which
is had and kept only through a capacity or skill to reproduce it.

This latter may now remain behind as merely bare, lifeless material, like a
residuum of the activity---as something now merely had and assumed because it was
once assumed, a \emph{standing result}, without its being any longer enlivened by
the living activity by which it was \emph{formed}. But it should not be
mechanized or crystallize itself in this way; rather, it should be kept
constantly ready, in the course of life's advance, to enter again under the
living influence of that determining First, and, under a constantly living and
organic rebirth, to receive transformations and new shapes. Only then is the
relation preserved as it is in itself and therefore also ought to be. As the
influences from above, from within, from that which is before us, direct and
command, so ought we to work in our lower and outer sphere.

\begin{center}Remark and Supplement.\end{center}
1.~In all cognizing by means of thinking (concept and Idea) this opposition and
doubleness is clear. Concepts, as such, are had and held fast only either by
means of the constructions whose norms they are, or by means of the judgments by
which their so-called content is determined. But the norm, or determining
fundamental thought, of a construction indeed enters and comes forward in it as
governing and ruling, yet is not itself constructed; it is therefore had only in
the living activity that proceeds from it. In judgments a concept is expressed,
in that its content is determined; but since this content is a sum of
determinations ascribed not to the concept itself but to its objects---yet in
accordance with the concept---the concept here too is had only as the determining
fundamental thought of a doing and determining, and is not embraced by it, but
hovers over it, or constitutes the inner first moving ruler within it, its inner
\textit{primus motor}. Which holds still more of every Idea. The concept can, to
be sure, in turn be regarded as an object that is itself determined; but this can
then happen only in accordance with another concept, which then constitutes the
ruling fundamental thought in the determinations. Such ruling fundamental
thoughts or concepts, which are had only in what is posited and carried out by
force of them, also underlie all inferences. In lower sensation, too, a
determining first enters into opposition to the representation thereby
determined. In striving to see and hear rightly and precisely, and to obtain a
complete representation of the object in which it can show itself in its full
distinctiveness, then---precisely because the need for one's own doing now comes
forward---that which is to be beheld will also show itself as that toward which
one must constantly look: that is, as that by which that doing is governed and
mastered, and which is obtained only by developing, in accordance with its own
determining influence, a complete image of it.

2.~The apprehending and cognizing individuality here shows itself as one which,
just as in its representations it forms a material in which it expresses its
cognizing and life, so too is itself formed into being that in which truth
expresses itself---quite as if the individuality were itself, in turn, the
substrate or material of the self-presenting truth, in which the latter seeks to
come to appearance and to take shape (cf. § 4, Remark 3). \emph{Thereby we grow},
on the one side, \emph{into that} whole in which a universal life expresses and
realizes itself, in that we are embraced, filled, and made alive by it as members
or organs; while on the other side a whole also grows and develops in us, in
which we ourselves realize ourselves as personal beings subsisting for ourselves.
We are ourselves assimilated by that more general life, in the very act of
seeking to assimilate to ourselves the true that presents itself to us, by
assimilating to ourselves its expression within us.

I refer here to the ingenious Franz Baader's treatise in Schelling and Markus's
\textit{Jahrbücher der Medicin}, third volume, first part, as well as to several
of his shorter philosophical-biological writings.

\begin{center}\large\S.~4.\end{center}
\markboth{\S~4}{\S~4}
In simple, pure, wholly immediate beholding, productivity is \emph{bound} and
\emph{lost}, or is so absorbed therein that it does not come into view on its
own, or come to consciousness, and stir according to its particular nature---just
as neither is the influence, with its determining and compelling power, felt and
brought to consciousness as such. The opposition between the cognized true and
the individuality in and for which it comes forward does not here show itself as
an opposition. This cognizing is also, as regards the manner of cognizing, the
most perfect; for the fact that the production comes entirely of itself proves
that truth works with complete power, and that it is also taken up fully.

But in a striving to cognize quite fully what presents itself immediately in this
beholding, and at the same time to cognize still other and more---because what is
thus cognized points to other and more---the \emph{object of cognition and the
cognizing personality}, which in pure beholding were undivided, step, by means of
a reflection, \emph{apart from one another} and against one another. Productivity
becomes unbound and shows itself as free, while the determining influence, on the
other side, shows itself as such, with its necessity. The individual finds itself
as one that must itself work, by way of deliberation and thinking, to reach the
desired cognition; and just for that reason productivity must show itself as a
free one, capable of presenting something other than what is found. But now,
therefore, the attachment too will be embraced by self-determination and will
occur in accordance with a willing. Truth now shows itself as an object that we
can come to possess only gradually, by effort, through a free activity that must
accord with truth's declarations---to which we must now attend closely, as they
gradually sound to us in the course of our inquiry. Whereby productivity, while
showing itself as free, is bound and mastered anew, so that it presents only in
the one true way---yet in such a way that it now occurs with self-determination,
and as one that could act otherwise, and therefore also by way of deliberation,
and such that a free choice is possible. But pure, wholly immediate
beholding---through a new, complete self-loss of the productivity that has come
to clarity and self-consciousness---is nonetheless that toward which one strives,
so that truth in its full fullness may enter into the cognizing personality.

\begin{center}Remark and Supplement.\end{center}
1.~The productivity that has become free can now also \emph{express itself in its
own} productions, independent of the given. But even therein it will always be
determined---since it must take place in definite ways---by certain norms or
concepts, which of themselves govern within it and are had directly in and
through the power with which they determine the self-activity. And if these
productions are to be more than a play with representations, then this freer
productivity---although its products are now not to be mere images of what is
found, but its own creations---must nonetheless, in the same way as in the
endeavor to cognize, be bound and mastered anew by a higher cognition, under
which the true shows itself to us in new forms and takes shape in us. So it is in
mathematics and in art.

In general, life comes forth just as fully in action and self-activity as in a
striving to cognize. In the former, free forth-working prevails, but also a
certain cognition, which underlies it as governing and moving First. But here too
the two opposites can, in the same way, be lost or merge into one, as in mere
beholding, so that neither is the self-activity particularly felt as free, nor
does the fundamental thought guiding it come to consciousness as such, since the
action happens entirely of itself: a self-activity that, as regards its manner of
working, is the most perfect---since therein the true and the being, the
essentiality that is to take shape in us in our acting too, shows itself as that
which works with complete, undivided, and unarrested power.

2.~So far as the free productivity comes to consciousness as such, cognition's
determining First too comes to consciousness as such; and vice versa. In a free
working at cognizing, as in an acting undertaken with self-conscious
self-determination, there can therefore, while \emph{something} comes
particularly \emph{into view}, be \emph{something else} deeper or more inward
which, with its determining influence, is lost or merges into the cognition or
the self-activity---in which respect the productivity, in another respect become
free, does not show itself as free. But a reflection can then afterward, in a
striving for deeper cognition, draw forth what was thus present hitherto but was
lost in the whole. ---Thus, for example, the principle of causality may at one
time underlie, as a determining fundamental thought, a striving to fathom certain
phenomena, without this principle being thought of; at another time reflection is
turned upon it; and the standpoint at which the cognizing striving stands in the
former case is quite other than in the latter.

3.~But why does the true, full essence of things not show itself at once to the
individual striving after cognition in such a way that all cognizing becomes a
pure, immediate beholding? It is natural that this question should arise here.
But its complete answer would require a general reflection on the ground of there
being a time and a temporality at all, and of life's coming forth at all in a
becoming [Vorden], under which it only gradually attains a definite
subsistence---something for whose full elucidation a fathoming of the innermost
meaning and essence of all human existence [Tilvær] would be needed. ---Only the
first basic features of such a consideration can find place here.

In the finite world the eternal and divine is to come forth, rule, and hold sway,
so as therein to reveal its infinite power and fullness: this is the meaning of
everything finite---that it is to be \emph{the Other} than the eternal, in which
the eternal is to come forth and reveal itself. But only in a consciousness and
spiritual life does the eternal and divine reveal itself \emph{as such}---in that
the being in which it lives comes to know of it, and moreover not only to bear it
within itself and be moved by it, but also itself to be able to behold and have
it as what it in its truth is. But the kingdom that is thus to come to the human
being is to be a \emph{kingdom of love}; for only in such a one does the true
fullness of the eternal come to show itself, and to show itself in its pure
spirituality, uniting a community of beings who are themselves spirits into a
whole that can now in truth be called the kingdom of God: the kingdom in which
God himself reveals himself in his divinity, in that he reveals himself in
eternal love. But a kingdom of spirit and love \emph{necessarily}
presupposes---and therefore necessarily calls forth---as the other in which it is
to be the ruling element, \emph{an intelligence and freedom}. If, on the part of
the beings to whom this kingdom is to come, a free and fully conscious devotion
did not come to take place---namely, a devotion of one's own mind, one's own
desire and impulse---then it would not be a kingdom of love that came forth. But
a personal life that is to be a free one---which now, out of pure love and
moreover with clear consciousness, is to give itself over to being enlivened and
filled by the divinity, and to move only in God---must itself make itself into
such a free life, in coming to clarity about the true meaning and essence of its
own existence, and likewise of the whole of existence. But then the opposition
must come forth fully, so that the reconciliation may become the true
reconciliation and unity in love and freedom. But if consciousness is thus to
come to self-consciousness and intelligence, telling itself what it is---if the
personality is to come to freedom under a self-determination in which it takes
itself as such and makes itself free---then it must come to a productivity in and
under which it objectifies itself for itself. It must come to express itself in a
sphere (that lower region mentioned above), which is its own, which it forms for
itself, and in which it becomes conscious of itself and of its freedom in a free
productivity. And thus life descends into a time and a temporality.

The personal being can, to be sure, no more make itself free (although we came to
use this expression) than it can give itself truth. It finds itself free as
regards possibility, but to have and enjoy freedom in deed (\textit{actu}) it
comes only through the liberating power of truth and love. But now it cannot be
liberated either, except \emph{in and through its own self-determination}---and
that a self-determination which it itself wills and feels as its own. That which
brings and gives the individual all reality, essentiality, and fullness can,
since it is a kingdom of love that is to form itself, give it the fullness only by
referring itself to \emph{something} that must now \emph{follow} in the individual
itself, as the individual's own self-determining: namely, its own self-surrender
and the acting that follows from it. The decision (that in which the whole working
concludes or ends) must always be the individual's own. Therefore, too---so that
we may still note this---the inner effects, in which the life that is to come to
reality operates or presses forward into consciousness, have the character of
incitement and summons: namely, a summons to seek what the individual can obtain
only when it desires and craves and seeks it, but which it can doubtless likewise
obtain only in that it is given to it outright. But of this again further on; see
§§ 19--21. In order that freedom may fully come to consciousness, and a complete
self-consciousness may set in---yet so that this self-consciousness also clearly
beholds what it is to strive to make itself into, so that now spirit and truth
may be had in spirit and truth---for this reason the latter opens itself to the
individual only gradually, under the individual's own striving, inquiry, and
labor.

\begin{center}\large\S.~5.\end{center}
\markboth{\S~5}{\S~5}
Insofar as the opposition between the object of cognition and the cognizing
personality comes forward in such a way that productivity, or representation,
shows itself as free---yet is now brought, by an incoming self-determination, to
present in one single way---cognition becomes a cognizing by force of grounds
[Grunde].

All cognizing is to be such a \emph{cognizing by force of grounds}.

For all cognizing is to come to complete clarity in and for the cognizing
personality and consciousness. But if it is to do so, then the productivity
required for it, and lying within it, must also come to clear consciousness; and
this it can do only in being carried out or undertaken by the individual with
self-consciousness and self-determination. Thereby it shows itself as free; but
if it is then, as that by which a cognition is to come about, bound to a single
determinate presenting, it must be bound to it in accordance with a certain
corresponding determining Something which, when it comes to consciousness as
such, shows itself as a ground.

Cognition now comes about as the result of an investigation and a deliberation.
Reflection and consideration call forth a \emph{study} that steps in between a
first immediate receiving and a final or concluding cognition. The mind, which
at first opened itself merely to receive and gave itself over purely to
contemplation and reception, is now set in motion in such a way that it is
itself, by its own self-conscious labor, to bring about the production in which
it is to possess truth; and it now seeks, by inquiry and grounding, a principle
or first mover for this production of its own. But the cognition that thereby
enters is nonetheless, when it now comes about, and as it now turns out, a new
receiving and beholding; for in the end we assume what, in accordance with the
determining grounds, we see that we must assume. This cognizing by force of
grounds is also to be a passage to a complete beholding, in which the full power
of truth works immediately and undividedly. But for the free personality, which
is to come to complete freedom and self-clarity, it shows itself as a necessary
passage.

\begin{center}Remark and Supplement.\end{center}
1.~The ground may, for that matter, well be merely a necessity accompanying pure
immediate apprehension and beholding, so that cognition takes place as one that
grounds itself; but the question of a ground has now been set in motion, and even
if what is assumed is assumed merely because it is immediately certain by itself,
this very immediate certainty [Vished] now shows itself as a ground, and
cognition now takes on a different character than it had before, when this
reflection had not yet entered.

This is especially to be remembered of the mathematical axioms. Sense-intuition,
too, no doubt presents itself as immediately grounding itself, and may also
remain with this stamp. But even then our representing is not seldom led to find
itself as a free reproducing---especially if the sensory shape is so manifold, or
the relations of its parts so intricate, that it can come to clarity only by
deliberation and consideration; and then what leads to fixing the representation
will also come, in opposition to the representation itself, to show itself as its
ground. This happens still more when it is noticed that things often cannot, or
should not, be represented as being just as the impression would immediately
suggest. But once such so-called sense-deceptions have shown themselves as
possible, doubt can now arise about the correctness of every immediate,
self-forming sense-representation, so that, even if it is directly assumed, this
now happens no longer merely on account of its own immediacy, but also on account
of a cognition that no so-called sense-deception is present---and the
consideration that teaches this now enters as a ground. Indeed, in the advance of
the endeavor and of cognition, the whole of sense-intuition finally shows itself
(to the philosophical eye) as one whose truth and validity can be doubted at all,
and must, moreover, first be warranted and established. And thus the representing
that at first formed itself entirely of itself, and seemed reliable enough in
itself, will show itself as still needing grounding---in general and in each and
every respect.

The intrinsically certain, from which this grounding is to proceed, or by which
all other cognition is to be warranted, is now sought on a wholly different side,
and cognition rises to a highest Idea, by force of which all else is now to be
beheld in its true and full light, in its meaning and truth. But even if this
highest Idea is, so far, immediate and resting in itself, we still come to it only
by a thinking that must, of course, throughout be guided and determined by
grounds---just as the very fact that all else can be clarified by that highest
Idea may be the ground on which its cognition rests. And thus we find, at the
highest point just as well, a cognizing by force of grounds, as we found the
lowest to be in need of grounding; and this because all cognizing comes to
actuality only through a being-determined to a certain representing, a
being-determined that, in the advance of the endeavor, must come to consciousness
as such. Everywhere the question is asked why we let ourselves be thus
determined, by what right this determining power is exercised over us.

2.~The perfect cognizing would be one that, independent of all deliberation and
grounding with all their unrest and striving movement, would rest wholly in
itself alone and be immediately clear and certain in itself, without any question
of its grounds arising, or any thought that something other than just what is
beheld could be the true. But such certainty and security---a cognition come to
rest and fullness in such a beholding---can rest only in the certainty of, indeed
in the pure immediate sense or feeling of, the fact that what is thus beheld
comes to us directly from the source of truth, is given and told to us by him
from whom all truth, like every good and perfect gift, comes. This is the
beholding that has been meant by the expression: to behold all things in God. It
is the cognizing by which we are to cognize even as we ourselves are
cognized---to know what we know because the spirit speaks it thus in us and gives
it to us. Then truth no longer stands over against us as the thing we seek;
rather, it dwells in us, in that the same that gives all things their
essentiality now also gives us cognition---pure, immediate, absolute, or in and
for itself, without its passing through any deliberation and reflection and
emerging as the result of such. A consciousness and assurance that everything we
cognize is given us by God---that our words are not ours, but, as we find them in
us, are worked in us by the divine spirit---can alone be that in which this
cognition can rest. Christ's cognizing we can conceive only in this way, never as
a cognizing by way of reflection and inference.

As a highest goal, as an ideal, such a cognition stands before us. Yet this
highest beholding cannot even be regarded as the goal that mere inquiry and
grounding, once it reached its completion, could attain. A completed thinking can
lead to a completed cognition by way of thinking---to a completed, wholly
certain, clear, and self-reposing knowledge according to, and by force of, the
highest grounds. But that highest beholding's pure immediacy can be reached only
in that the highest true---the eternal source of truth, the spirit---becomes
living in us and enlivens us in a wholly different (a differently real) way than
the way in which truth, as highest Idea, animates and enlivens a thinking and a
grounding.

For the forward-striving human being, however, \emph{the way of thinking} is the
one along which he is to advance, in order to press more and more toward the full
cognition of truth. It is a more earthly way, on which one attains piece by
piece. It belongs to the sphere in which, within a finitude, under its own
striving, a personality is to develop and be raised more and more to clarity and
independence---in a life in which it more and more objectifies itself for itself
and establishes itself in an existence of its own, yet only in order to ripen to
the point where the spirit can now, in its fullness, enter into it as into an
individuality now fully developed as its organ. Whereby then, when this point is
reached, fullness will enter in place of the piecework; and now thinking---which
surely is not to cease---and the fathoming of things---which is not to fall
away---will themselves be illuminated by the higher light, enter into the higher
peace, and, while an infinite fullness stands before them for contemplation and
fathoming, will nonetheless themselves walk the secure course that no longer has
constantly to ask itself about the grounds why it assumes what it beholds, why it
cognizes what it finds in its grounding, but has in this respect cast out ``every
witness of human neediness.''

Until it comes to that, all cognition is burdened with \emph{uncertainty}. The
immediately given, which is directly beheld, cannot satisfy us, because it is not
the true essence on whose cognition everything depends, and because we have of
this true essence a fundamental consciousness that drives us to seek something
other than what is directly beheld. It itself now points us to something else and
higher; but this else and higher we come to only precisely by being thus shown
and led to it---that is, in a mediated way---and this mediateness remains, so long
as our cognizing is a cognizing by way of thinking, always attached to it for us.

But what can here already lift us above this mediateness is the firm
\emph{conviction} [Overbeviisning], forming itself, of its reality. This springs
up in the course of our inquiry and grounding, as we are seized by the truth of
what is cognized, feel ourselves carried away by it. In the course of our
comparing and relating, deliberating and considering, it lifts itself more and
more forward and becomes---in the subjective, individual, and personal respect
(what concept and insight are in the objective respect)---the moving soul of all
thinking and determining, from which it proceeds and to which it finally returns.
And thus a \emph{faith} [Tro]---overcoming and dissolving all deficiency in
insight and in the developments and fixings of thinking---namely, a \emph{faith in
the reality of the fundamental cognition [Grunderkjendelse]}---takes for us the
place of that perfect cognition. As faith it is---though awakened and nourished by
a grounding, and constantly giving birth to a closer deliberation and grounding of
itself---nonetheless not a cognizing by, and by force of, the development of
grounds alone, but a more immediate appearance of the power of truth, which,
bearing witness to itself, seizes and irresistibly compels to cognition---yet not
because there is immediate beholding, but because everything so points and presses
us to this cognition that it shows itself as the center on which everything
depends and in which everything rests, and without which we cannot subsist or
manage at all. As our life in general, so too our life in cognition must seek its
last support and rest in an infinite trust and confidence, and, as the life of a
finite being, can in the end rest in nothing else. Yet that here too, as
everywhere, divine grace is in the end the one thing on which everything depends
and in which everything rests---this we shall come to set before ourselves more
closely in what follows. But because we receive only by grace, that too is why we
can have only by faith. Everything is in the end given, and must be taken as it is
given, in faith toward the Giver. But precisely that which makes cognition not
become a total beholding now also opens for this infinite trust the room in which
it can show itself as infinite trust---so that the very fact that we do not have
the fullness brings it about that attachment, devotion, and love come all the
more into view.

And that, then, is what all the life in us aims at, and moreover also all the life
that labors to reach a perfect cognition: that we, the more we exert ourselves to
have something of our own and to possess something in ourselves, should the more
come to feel that we can do nothing and have nothing except what is directly and
outright given us; and that, the more we rise in independence and intelligence and
command---as regards the manifold and its relative determination---the more we
should recognize that, in respect of the One thing that is needful and on which
everything depends, we must give ourselves entirely into God's keeping, who alone
works in us both to cognize and to carry out.

3.~But it is only of the fundamental cognition (that on which everything depends,
and which so far is the One thing needful) that this holds---that it rests, in the
end, in a faith and trust. Everything else, which is illuminated by this
fundamental cognition and is to be determined in accordance with it, must be
beheld as given with it---all in its connection therewith and its consequence
therefrom. So that, even if on the one side a faith and trust becomes that in
which all conviction rests and cannot but rest, cognition is nonetheless not on
that account to cease, on the other side, being a cognition by force of grounds,
and of an \emph{insight} arisen through the cognition of grounds. And just as
everything cognized in accordance with the fundamental cognition must be cognized
so, because we clearly see that the fundamental cognition necessarily entails it,
and can follow the whole connection from the single point under consideration up
to the highest, so too the fundamental cognition itself must precisely be clearly
seen as that which bears the whole connection, and to which this connection
necessarily leads. Under which the course will be this: that what was perhaps at
first cognized by a certain feeling of truth---indeed, by a relation to the
subjectively human---comes more and more to be evident to thinking in and for
itself and in an objective way, until at last a point will be reached where the
highest, to which all else points, is evident not merely because it is clearly
seen that all else points to it, but because it seems immediately \emph{evident in
itself alone}. But this consideration---in which we meet the profound thought
underlying the famous so-called ontological proof of God's existence---leads us
over to our next principal proposition. ---God bears witness to himself in
consciousnesses, and no cognition that God is would be possible if God did not
thus bear witness to himself within us, that is, in our gradually forming
spiritual life, beholding, and contemplation; indeed, this fundamental
consciousness of God (and of one's own innermost communion with God) is what
first makes the human being human. Herein, then, is immediate faith. But should
not the same cognition also be able to become just as immediately evident for
speculation, so that this evidence too, as evidence, rested only in itself? By
which, however---since this evidence too would then have to be given to us, in
that it was evident precisely by itself---faith would not cease to be that in
which all cognition would in the end rest. Insight and self-evidence, or evidence,
are themselves \emph{worked} in us as well; but where there is question of a life
working itself forward in consciousness in a real way, and what comes forth is
seen as this life's inner operations, there is grace and governance---and there,
moreover, faith, attachment, and trust must show themselves. The purely
intelligent is itself something that likewise comes only through real influences
from above, and must therefore be received as a gift---and that with trust toward
the Giver. In the cognizing that we conceive as Christ's cognizing, an immediate
self-evidence becomes one with the assurance that trust gives; it is faith and
beholding in one.

\begin{center}\large\S.~6.\end{center}
\markboth{\S~6}{\S~6}
% [text to be added: pp. 48--54 (\S 6, main text)]

\begin{center}Remark and Supplement.\end{center}
% [text to be added: pp. 48--54 (\S 6, Remark and Supplement)]

\begin{center}\large\S.~7.\end{center}
\markboth{\S~7}{\S~7}
% [text to be added: pp. 55--71 (\S 7, main text)]

\begin{center}Remark and Supplement.\end{center}
% [text to be added: pp. 55--71 (\S 7, Remark and Supplement)]

\begin{center}\large\S.~8.\end{center}
\markboth{\S~8}{\S~8}
% [text to be added: pp. 72--87 (\S 8, main text)]

\begin{center}Remark and Supplement.\end{center}
% [text to be added: pp. 72--87 (\S 8, Remark and Supplement)]

\begin{center}\large\S.~9.\end{center}
\markboth{\S~9}{\S~9}
% [text to be added: pp. 88--103 (\S 9, main text)]

\begin{center}Remark and Supplement.\end{center}
% [text to be added: pp. 88--103 (\S 9, Remark and Supplement)]

\begin{center}\large\S.~10.\end{center}
\markboth{\S~10}{\S~10}
% [text to be added: pp. 104--111 (\S 10, main text)]

\begin{center}Remark and Supplement.\end{center}
% [text to be added: pp. 104--111 (\S 10, Remark and Supplement)]

\begin{center}\large\S.~11.\end{center}
\markboth{\S~11}{\S~11}
% [text to be added: pp. 112--115 (\S 11, main text)]

\begin{center}Remark and Supplement.\end{center}
% [text to be added: pp. 112--115 (\S 11, Remark and Supplement)]

\begin{center}\large\S.~12.\end{center}
\markboth{\S~12}{\S~12}
% [text to be added: pp. 116--119 (\S 12, main text)]

\begin{center}Remark and Supplement.\end{center}
% [text to be added: pp. 116--119 (\S 12, Remark and Supplement)]

\begin{center}\large\S.~13.\end{center}
\markboth{\S~13}{\S~13}
% [text to be added: pp. 120--127 (\S 13, main text)]

\begin{center}Remark and Supplement.\end{center}
% [text to be added: pp. 120--127 (\S 13, Remark and Supplement)]

\begin{center}\large\S.~14.\end{center}
\markboth{\S~14}{\S~14}
% [text to be added: pp. 128--132 (\S 14, main text)]

\begin{center}Remark and Supplement.\end{center}
% [text to be added: pp. 128--132 (\S 14, Remark and Supplement)]

\begin{center}\large\S.~15.\end{center}
\markboth{\S~15}{\S~15}
% [text to be added: pp. 133--145 (\S 15, main text)]

\begin{center}Remark and Supplement.\end{center}
% [text to be added: pp. 133--145 (\S 15, Remark and Supplement)]

\begin{center}\large\S.~16.\end{center}
\markboth{\S~16}{\S~16}
% [text to be added: pp. 146--152 (\S 16, main text)]

\begin{center}Remark and Supplement.\end{center}
% [text to be added: pp. 146--152 (\S 16, Remark and Supplement)]

\begin{center}\large\S.~17.\end{center}
\markboth{\S~17}{\S~17}
% [text to be added: pp. 153--158 (\S 17, main text)]

\begin{center}Remark and Supplement.\end{center}
% [text to be added: pp. 153--158 (\S 17, Remark and Supplement)]

\begin{center}\large\S.~18.\end{center}
\markboth{\S~18}{\S~18}
% [text to be added: pp. 159--159 (\S 18, main text)]

\begin{center}\large\S.~19.\end{center}
\markboth{\S~19}{\S~19}
% [text to be added: pp. 160--161 (\S 19, main text)]

\begin{center}\large\S.~20.\end{center}
\markboth{\S~20}{\S~20}
% [text to be added: pp. 162--164 (\S 20, main text)]

\begin{center}Remark and Supplement.\end{center}
% [text to be added: pp. 162--164 (\S 20, Remark and Supplement)]

\begin{center}\large\S.~21.\end{center}
\markboth{\S~21}{\S~21}
% [text to be added: pp. 165--193 (\S 21, main text)]

\begin{center}Remark and Supplement.\end{center}
% [text to be added: pp. 165--193 (\S 21, Remark and Supplement)]

\begin{center}\large\S.~22.\end{center}
\markboth{\S~22}{\S~22}
% [text to be added: pp. 194--204 (\S 22, main text)]

\begin{center}Remark and Supplement.\end{center}
% [text to be added: pp. 194--204 (\S 22, Remark and Supplement)]

\begin{center}\large Conclusion.\end{center}
% [text to be added: Slutning (Conclusion), pp. 203--204]

% [Rettelser / errata leaf, printed p.205: corrections to the DANISH text
%  (page/line orthography fixes). Meaningless in English translation.
%  DECISION DEFERRED (see TRANSLATION-NOTES): omit, or keep as a short
%  translator's note recording that the Danish had an errata leaf.]

\end{document}
